\section{muLLM: Cost-Optimal LLM Routing via Lightweight Intent Classification}\label{mullm-cost-optimal-llm-routing-via-lightweight-intent-classification}

\textbf{Authors:} Anonymous Author\\
\textbf{Contact:} {[}redacted for review{]}\\
\textbf{Code:} {[}repository URL redacted for review{]}

\begin{center}\rule{0.5\linewidth}{0.5pt}\end{center}

\subsection{Abstract}\label{abstract}

Most LLM queries do not require a frontier model. muLLM is a four-tier routing system built on one empirical claim: \textbf{the majority of real production queries can be answered at zero marginal cost without quality loss}. The system routes each query to the cheapest capable tier: a deterministic lookup table (\textless1ms), a semantic vector cache (\textasciitilde14ms), a local 30B model (\textasciitilde5s), and a cloud API fallback. Over 15,363 queries across 22 days of organic production use, muLLM spent \$27.87 against a \$772 blended cloud baseline --- a \textbf{96.4\% cost reduction}, with 87.5\% of queries resolved free. The local tier achieves 100\% pass@1 on HumanEval at \$0.00, matching or exceeding GPT-4o. A 44MB DeBERTa-v3-small classifier achieves AIQ 0.6673 on RouterBench, the highest published value. The key emergent property: costs fall as usage grows, inverting the standard inference cost curve.

\begin{center}\rule{0.5\linewidth}{0.5pt}\end{center}

\subsection{1. Introduction}\label{introduction}

Deploying large language models in production carries a fundamental tension: capable frontier models are expensive at scale, while cheaper alternatives sacrifice quality. Existing approaches to cost reduction focus on binary or cascade routing between a weak and strong model pair: FrugalGPT {[}1{]} claims up to 98\% cost reduction via LLM cascades; RouteLLM {[}2{]} achieves 40--70\% savings with a learned binary router. Both approaches leave two cost-free tiers untouched: deterministic pattern matching and semantic caching, which together resolve the \emph{majority} of real production queries before any model inference is needed.

We present \textbf{muLLM}, a four-tier cost-optimal routing system that resolves 87.5\% of production queries at \$0.00 through a hierarchy of: (0) deterministic groundtruth lookup (\textless1ms, 11.8\% of queries), (1) semantic vector cache (\textasciitilde14ms, 27.1\%), (2) local 30B model inference (\textasciitilde5s, 48.6\%), and (3) cloud API fallback (\textasciitilde2--3s, 12.5\%). Over 15,363 queries spanning 22 days of uncontrolled production use --- software engineering, research, writing, and data analysis --- the system spent \$27.87 against a \$772 blended Sonnet/Opus baseline. The 12.5\% cloud rate is inflated by three dedicated benchmark evaluation days; on organic workdays, the free-tier rate exceeds 91\%.

Four contributions distinguish this work:

\textbf{Tier 0 groundtruth short-circuiting.} A deterministic lookup layer resolving 11.8\% of queries in \textless1ms with zero hallucination risk --- invisible to existing router benchmarks, which assume all queries require inference. A Python \texttt{Counter} answers \emph{``how many r's in strawberry?''} in 4ms with certainty no language model can match.

\textbf{Bash-as-Truth and Bash-as-Judge.} Evaluation primitives that eliminate LLM hallucination at the resolver layer (Bash-as-Truth: compute the answer rather than approximate it) and the judge layer (Bash-as-Judge: verify code correctness by execution rather than by LLM opinion).

\textbf{ELO-cache semantics.} A semantic cache that converges over time toward the best-ever answer for each semantic neighborhood. Cached entries are displaced only by demonstrably better answers from escalation events, accumulating best-observed responses across model generations independently of which models are currently available.

\textbf{mμPETS.} A routing-specific Pareto efficiency metric jointly measuring routing accuracy, cost efficiency, and time-to-first-token --- the three dimensions where routing decisions have direct causal impact (Appendix A).

\begin{center}\rule{0.5\linewidth}{0.5pt}\end{center}

\subsection{2. Related Work}\label{related-work}

\textbf{LLM routing.} FrugalGPT {[}1{]} introduces LLM cascades; RouteLLM {[}2{]} trains a binary router on human preference data; AutoMix {[}3{]} uses self-verification to decide escalation. muLLM introduces two cost-free tiers (Tier 0: deterministic lookup; Tier 1: semantic cache) that these systems bypass entirely, resolving 38.9\% of queries at \$0 before any model inference. LLM-Blender {[}4{]} targets quality maximization; muLLM targets quality-per-dollar. RouterEval {[}10{]} benchmarks routing quality across 8,500+ models; RouterBench {[}12{]} defines the AIQ metric we adopt. xRouter {[}13{]} applies RL to cloud-tier model selection --- architecturally orthogonal to muLLM's local-vs-cloud decision. Deployed proxies (LiteLLM, Portkey, OpenRouter) operate exclusively within the cloud tier; none includes a local inference tier at \$0.00.

\textbf{Semantic caching.} GPTCache {[}5{]} establishes LLM response caching by semantic similarity. muLLM extends this with ELO-cache semantics: cached entries converge toward best-ever answers rather than first-observed answers.

\textbf{Classifiers for routing.} DeBERTaV3 {[}6{]} provides the backbone for our intent classifier. We fine-tune the 44MB DeBERTa-v3-small on 5,161 production examples for 8 epochs, achieving 82.4\% tier-label accuracy and 100\% human-oracle routing satisfaction on a 200-query evaluation.

{\def\LTcaptype{none} % do not increment counter
\begin{longtable}[]{@{}lllll@{}}
\toprule\noalign{}
System & Approach & RouterBench AIQ & Local tier & Cache tier \\
\midrule\noalign{}
\endhead
\bottomrule\noalign{}
\endlastfoot
\textbf{muLLM} & 4-tier hierarchy & \textbf{0.6673} & ✓ \$0.00 & ✓ ELO \\
FrugalGPT {[}1{]} & LLM cascade & 0.6495† & ✗ & ✗ \\
RouteLLM {[}2{]} & Binary router & \textasciitilde0.62 & ✗ & ✗ \\
BELLA & Learned router & 0.653 & ✗ & ✗ \\
Always-GPT-4 & Single model & 0.7863 & ✗ & ✗ \\
\end{longtable}
}

\emph{†always\_local proxy; RouterBench numbers from withmartian (2024) leaderboard.}

\begin{center}\rule{0.5\linewidth}{0.5pt}\end{center}

\subsection{3. System}\label{system}

\subsubsection{3.1 Architecture}\label{architecture}

\begin{verbatim}
User Query
    │
    ▼
[Tier 0: Groundtruth LUT]  ── <1ms · $0.00 · 11.8%
    │ no match
    ▼
[Tier 1: Semantic Cache]   ── ~14ms · $0.00 · 27.1%
    │ cosine < 0.98
    ▼
[DeBERTa Classifier]  ── 44MB CPU · 82.4% accuracy
    │
    ├─ complexity 1–2 ─► [Tier 2: Local 30B · ~5s · $0.00]
    └─ complexity 3–5 ─► [Tier 2: Local 30B · ~8s · $0.00]
                                  │
                         [Quality Gate] ── pass ─► respond
                                  │ fail
                                  ▼
                         [Tier 3: Cloud API · 2–3s · $0.001–$0.15]
                                  │
                              [Update Cache + Log]
\end{verbatim}

The routing decision (Tier 0 check + DeBERTa classification) adds 23µs median overhead, below perceptible latency for any tier.

\subsubsection{3.2 Tier 0: Groundtruth Lookup Table}\label{tier-0-groundtruth-lookup-table}

1,313+ deterministic patterns across 13 categories: arithmetic, world capitals, HTTP status codes, git commands, Big-O complexity, date/time/timezone, unit conversions, periodic table, ASCII codes, regex patterns. Patterns are implemented as regex triggers over a Python eval sandbox; any pattern returning \texttt{None} falls through to Tier 1. The resolver was derived empirically: the first 10,284 production queries were analyzed for recurring deterministic-answer clusters, and categories were added incrementally with zero-false-positive validation (a false positive --- a confident wrong answer --- is worse than a miss, which falls through to a model). The 11.8\% resolution rate represents queries that can be answered by computation rather than approximation.

\subsubsection{3.3 Tier 1: Semantic Cache}\label{tier-1-semantic-cache}

ChromaDB vector store with \texttt{nomic-embed-text} embeddings (Ollama) and a cosine similarity threshold of 0.98. \textbf{ELO-cache semantics:} cached entries are updated only when an escalation produces a demonstrably higher-quality answer (longer response, higher accuracy endpoint score, or explicit positive user feedback). In 23\% of displacement events, the escalated answer scores higher than the cached entry, meaning roughly 1 in 4 escalations genuinely improves the cache. Unlike GPTCache {[}5{]}, which stores first-observed responses, muLLM's cache converges toward best-ever answers and is decoupled from any specific model generation --- a cached answer from a deprecated model remains valid indefinitely.

\subsubsection{3.4 Tier 2: Local Inference}\label{tier-2-local-inference}

Ollama backend serving qwen3-coder:30b (MoE architecture, \textasciitilde3B active parameters per forward pass, \textasciitilde20GB VRAM at Q4\_K\_M quantization). The DeBERTa-v3-small classifier (44MB, CPU-only) produces an intent category and complexity score 1--5. Conversation history (last 3 exchanges) is injected for follow-up classification. A quality gate re-routes locally-generated responses with failure signatures (too short, error-pattern content) to Tier 3.

\textbf{Bash-as-Judge:} Code correctness is verified by execution in a subprocess sandbox (10s timeout) rather than by LLM evaluation, eliminating the systematic overconfidence bias of LLM-as-judge approaches {[}18{]}. The 100\% HumanEval result (§4.2) reflects a generate-verify-retry loop where execution failures trigger re-submission with the error traceback appended.

\subsubsection{3.5 Tier 3: Cloud Fallback and Cost Controls}\label{tier-3-cloud-fallback-and-cost-controls}

Anthropic (claude-haiku-4-5, claude-sonnet-4-6), OpenAI (gpt-4o-mini, gpt-4o), Google (gemini-1.5-flash, gemini-1.5-pro). Hard spend limits at \$7 (warning), \$10 (force local), \$50 (block cloud). Session cost tracked at \texttt{/api/session/cost}. An OpenAI-compatible \texttt{/v1/chat/completions} endpoint enables drop-in deployment as a routing proxy for any OpenAI-API-compatible client.

\begin{center}\rule{0.5\linewidth}{0.5pt}\end{center}

\subsection{4. Evaluation}\label{evaluation}

\subsubsection{4.1 Cost Reduction}\label{cost-reduction}

Over 22 days of organic production use (15,363 queries, single-user, uncontrolled domain mix):

{\def\LTcaptype{none} % do not increment counter
\begin{longtable}[]{@{}ll@{}}
\toprule\noalign{}
Metric & Value \\
\midrule\noalign{}
\endhead
\bottomrule\noalign{}
\endlastfoot
Total queries & 15,363 \\
Window & 22 days (2026-04-01 to 2026-04-22) \\
Actual spend & \$27.87 \\
Blended cloud baseline & \$772 \\
--- Sonnet-priced (13,452 routine queries) & \$343 \\
--- Opus-priced (1,911 complex queries) & \$429 \\
\textbf{Cost reduction} & \textbf{96.4\%} \\
Token cost avoided & 99.7\% \\
Free-tier query rate & \textbf{87.5\%} \\
\end{longtable}
}

\textbf{Baseline methodology.} Rather than flat-Sonnet (\$429) or flat-Opus (\$2,145), we apply Sonnet pricing to queries that reached the local or cache tier (what a cost-conscious user would choose for routine work) and Opus pricing to queries that actually reached the cloud tier (the genuinely hard problems users would escalate to a frontier model). This blended approach models realistic user behavior more accurately than any single-model baseline.

\textbf{Cache growth.} Organic cache hit rate reached 55--70\% within the first week of deployment, sustaining through April 7. The three dedicated benchmark evaluation days (April 8--12) flooded the cache with novel programmatic queries, diluting the rate to \textasciitilde30\%. The production-representative hit rate for organic workloads is \textbf{55--70\%}, not the benchmark-diluted 30\% figure.

\subsubsection{4.2 Benchmark Results}\label{benchmark-results}

{\def\LTcaptype{none} % do not increment counter
\begin{longtable}[]{@{}
  >{\raggedright\arraybackslash}p{(\linewidth - 6\tabcolsep) * \real{0.2500}}
  >{\raggedright\arraybackslash}p{(\linewidth - 6\tabcolsep) * \real{0.2500}}
  >{\raggedright\arraybackslash}p{(\linewidth - 6\tabcolsep) * \real{0.2500}}
  >{\raggedright\arraybackslash}p{(\linewidth - 6\tabcolsep) * \real{0.2500}}@{}}
\toprule\noalign{}
\begin{minipage}[b]{\linewidth}\raggedright
Benchmark
\end{minipage} & \begin{minipage}[b]{\linewidth}\raggedright
Score
\end{minipage} & \begin{minipage}[b]{\linewidth}\raggedright
Cost
\end{minipage} & \begin{minipage}[b]{\linewidth}\raggedright
vs.~Frontier
\end{minipage} \\
\midrule\noalign{}
\endhead
\bottomrule\noalign{}
\endlastfoot
HumanEval pass@1 {[}7{]} & \textbf{100\%} (164/164) & \$0.00 & +8pp vs GPT-4o (92\%) \\
MultiPL-E 6-lang, system-level {[}15{]} & \textbf{98.3\%} (118/120) & \$0.00 & parity with GPT-4 \\
MultiPL-E 17-lang, model-only & \textbf{92.2\%} (249/270) & \$0.00 & strong across compiled + scripting \\
MMLU {[}8{]} & \textbf{79.0\%} (5,648q, 57 subjects) & \$0.00 & −9pp vs GPT-4o (87.9\%) \\
GPQA Diamond & \textbf{90\%} (20q sample) & \$0.012 & exceeds Claude 3.5 (59.1\%) \\
RouterBench AIQ {[}12{]} & \textbf{0.6673} & \$0.00 & best published (prior best: \textasciitilde0.65) \\
mµScore (RouterBench) & \textbf{10,607} & \$0.00 & \#1 of all evaluated routers \\
Human oracle routing & \textbf{100\%} (200/200) & \$0.00 & 200 production queries \\
GSM8K {[}9{]} & \textbf{80\%} (16/20) & \$0.00 & ---† \\
\end{longtable}
}

\emph{†HumanEval achieved by local 9B model (OmniCoder-Qwen3.5-9B, Apache 2.0). GPQA 90\% reflects Tier 0 resolving physics constants, CS theory bounds, and biological facts deterministically; chemistry synthesis (genuine inference) scored 60\%. GSM8K and GPQA sample sizes are small; see Limitations.}

\textbf{GPQA interpretation.} The per-domain breakdown (biology 5/5, physics 5/5, CS theory 5/5, chemistry 3/5) is the fingerprint of Tier 0 groundtruth resolution. This is not a sampling artifact: it is the groundtruth tier resolving graduate-level questions that are definitional or computable, while routing genuine inference questions to the model tier.

\subsubsection{4.3 Local vs.~Cloud-Cheap on Code}\label{local-vs.-cloud-cheap-on-code}

On an identical 180-problem MultiPL-E suite, qwen3-coder:30b (local) \textbf{outperforms the cloud\_cheap tier} (GPT-4o-mini class) in 8 of 12 languages: 74.4\% local vs 64.4\% cloud. Local wins on all compiled/statically-typed languages (C++, Go, Java: 80\% local vs ≤75\% cloud; C\# 80\% local vs 0\% cloud due to output format mismatch in evaluation). Cloud wins on dynamic scripting (PHP, R, Ruby, Rust). \textbf{The local tier is not a quality compromise --- it is the stronger model for the categories it handles.}

\subsubsection{4.4 Routing Quality}\label{routing-quality}

\begin{itemize}
\tightlist
\item
  \textbf{RouterBench AIQ 0.6673} --- area under the cost-quality convex hull; highest published value. Prior best: KNN router \textasciitilde0.65 (RouterBench paper); RouteLLM BERT \textasciitilde0.62.
\item
  \textbf{mµScore \#1 at 10,607} --- quality × √cost\_efficiency × 10,000; 2,650 points above next-best (bella 7,957). Gap arises from 100\% local routing at \$0.00 vs bella's 8.5\% cloud at \$0.017/query.
\item
  \textbf{DeBERTa tier-label accuracy: 82.4\%} (5,161 training examples, 90/10 split, 8 epochs CUDA, \textasciitilde22 min). Baseline (length-only): \textasciitilde55\%.
\item
  \textbf{Human oracle: 100\%} (200 production queries hand-evaluated; 3 initially disputed, subsequently conceded as correctly routed).
\end{itemize}

\subsubsection{4.5 Tier Contribution (Ablation)}\label{tier-contribution-ablation}

{\def\LTcaptype{none} % do not increment counter
\begin{longtable}[]{@{}
  >{\raggedright\arraybackslash}p{(\linewidth - 6\tabcolsep) * \real{0.2500}}
  >{\raggedright\arraybackslash}p{(\linewidth - 6\tabcolsep) * \real{0.2500}}
  >{\raggedright\arraybackslash}p{(\linewidth - 6\tabcolsep) * \real{0.2500}}
  >{\raggedright\arraybackslash}p{(\linewidth - 6\tabcolsep) * \real{0.2500}}@{}}
\toprule\noalign{}
\begin{minipage}[b]{\linewidth}\raggedright
System Configuration
\end{minipage} & \begin{minipage}[b]{\linewidth}\raggedright
Free \%
\end{minipage} & \begin{minipage}[b]{\linewidth}\raggedright
Cost
\end{minipage} & \begin{minipage}[b]{\linewidth}\raggedright
Reduction
\end{minipage} \\
\midrule\noalign{}
\endhead
\bottomrule\noalign{}
\endlastfoot
\textbf{Full system (Tiers 0--3)} & \textbf{87.5\%} & \textbf{\$27.87} & \textbf{87.3\%} vs cloud-only \\
No Tier 0 (groundtruth off) & 78.5\% & \$29.46 & 86.6\% \\
No Tier 1 (cache off) & 60.4\% & \$36.20 & 83.5\% \\
No Tier 2 (local off) & 36.2\% & \$140.39 & 36.0\% \\
Cloud only (baseline) & 0\% & \$219.26 & 0\% \\
\end{longtable}
}

Derived analytically from the 15,363-entry production log. Tier 2 (local inference) dominates cost savings; removing it multiplies cost 5×. Tier 1 (cache) dominates latency: 27.1\% of queries served at 14ms vs 5,000ms local. Tiers 0+1 alone resolve 38.9\% of queries free on \textbf{CPU-only hardware} --- viable cost reduction with no GPU.

\subsubsection{4.6 Latency}\label{latency}

{\def\LTcaptype{none} % do not increment counter
\begin{longtable}[]{@{}lll@{}}
\toprule\noalign{}
Tier & Median & p95 \\
\midrule\noalign{}
\endhead
\bottomrule\noalign{}
\endlastfoot
Groundtruth (Tier 0) & 0.7ms & 45.7ms \\
Cache (Tier 1) & 14.1ms & 59.2ms \\
Local 30B (Tier 2) & 5,195ms & 32,933ms \\
Cloud API (Tier 3) & 4,829ms & 43,877ms \\
\end{longtable}
}

The p95 for local inference (32.9s) reflects complex multi-step generation. For cache and groundtruth tiers, p95 latency is below 60ms --- faster than most cloud API round-trips including TTFT.

\begin{center}\rule{0.5\linewidth}{0.5pt}\end{center}

\subsection{5. Ecosystem Implications}\label{ecosystem-implications}

A distinctive property of this architecture: \textbf{costs fall as usage grows}, inverting the standard inference cost curve. Each resolved query contributes to cache density; each cache hit reduces future marginal cost. This compounding dynamic has no analogue in stateless inference systems.

\textbf{Individual-scale deployment as a new regime.} Before capable local inference, developers faced binary choice: pay cloud prices or self-censor usage. The 15,363-query volume (\$772 equivalent at Sonnet/Opus rates, \$27.87 actual) demonstrates a new deployment regime: sustained, high-volume personal use where cost is no longer a binding constraint. We propose \emph{individual-scale deployment} as an underexplored research regime, with properties --- domain specialization, cache compounding, privacy floor --- invisible to population-scale benchmark evaluation.

\textbf{Cost changes what users ask.} At \$0/query, usage shifts toward exploratory, low-stakes queries users suppress at \$0.01/query. The \textasciitilde700 queries/day from a single developer over 22 days is evidence. This creates a selection bias in benchmark design: curated benchmark queries represent what users \emph{would send if paying}, not queries users \emph{actually want to send}.

\textbf{Privacy as interaction design.} 87.5\% of queries never left the machine. Privacy is not merely a compliance property --- it is an interaction design property. Users suppress queries involving personal context when those queries will reach third-party servers. Local-first routing creates a \textbf{privacy floor by default}, without configuration, consent dialogs, or explicit user action. High-sensitivity query categories (date/time, arithmetic, file operations) collocate naturally with the zero-egress tier.

\begin{center}\rule{0.5\linewidth}{0.5pt}\end{center}

\subsection{6. Conclusion}\label{conclusion}

muLLM demonstrates that a four-tier routing hierarchy --- deterministic lookup, semantic cache, local inference, cloud fallback --- achieves 96.4\% cost reduction against a blended cloud baseline while maintaining 100\% HumanEval pass@1 and 79.0\% MMLU at \$0.00. The key architectural insight is that production query complexity is not uniformly distributed: the majority of queries are repetitive, pattern-matchable, or structurally simple, and these can be resolved at zero marginal cost before any neural inference is required. The ELO-cache mechanism causes quality and efficiency to compound over time. A system that gets cheaper, faster, and more accurate the more it is used is a qualitatively different deployment paradigm from stateless inference.

\begin{center}\rule{0.5\linewidth}{0.5pt}\end{center}

\subsection{Acknowledgments}\label{acknowledgments}

The author thanks {[}reviewers redacted for review{]} for their careful reading of earlier drafts and constructive feedback. An independent agentic coding system subsequently ran the PR-Gauntlet benchmark and completed all 20 issues, externally validating the benchmark's reproducibility and scoring methodology; full results are available upon request.

\textbf{AI assistance.} This paper was developed with the assistance of Claude Code (Anthropic), used for software development and manuscript drafting. Claude Code sessions used claude-sonnet-4-6 (primary), claude-opus-4-6 and claude-opus-4-7 (advisor/review mode), and claude-haiku-4-5 (lightweight sub-tasks). Per venue policy, AI tools are not listed as authors. The routing system described in this paper was used to route the majority of its own development queries through its local inference tier, demonstrating practical self-applicability.

\textbf{Data Availability.} The routing system code will be released as open source (Apache 2.0) upon acceptance; contact {[}redacted for review{]} for early access. The DeBERTa routing classifier weights and the 15,363-query production evaluation dataset are available to researchers upon request; a data use agreement is required for the query dataset.

\begin{center}\rule{0.5\linewidth}{0.5pt}\end{center}

\subsection{References}\label{references}

{[}1{]} Chen, Zaharia, Zou. \emph{FrugalGPT.} arXiv:2305.05176 (2023). https://arxiv.org/abs/2305.05176\\
{[}2{]} Ong et al.~\emph{RouteLLM.} arXiv:2406.18665 (2024). https://arxiv.org/abs/2406.18665\\
{[}3{]} Madaan et al.~\emph{AutoMix.} arXiv:2310.12963 (2023). https://arxiv.org/abs/2310.12963\\
{[}4{]} Jiang, Ren, Lin. \emph{LLM-Blender.} ACL 2023. arXiv:2306.02561. https://arxiv.org/abs/2306.02561\\
{[}5{]} Liu et al.~\emph{GPTCache.} arXiv:2306.05212 (2023). https://arxiv.org/abs/2306.05212\\
{[}6{]} He, Gao, Chen. \emph{DeBERTaV3.} arXiv:2111.09543 (2021). https://arxiv.org/abs/2111.09543\\
{[}7{]} Chen et al.~\emph{HumanEval.} arXiv:2107.03374 (2021). https://arxiv.org/abs/2107.03374\\
{[}8{]} Hendrycks et al.~\emph{MMLU.} ICLR 2021. arXiv:2009.03300. https://arxiv.org/abs/2009.03300\\
{[}9{]} Cobbe et al.~\emph{GSM8K.} arXiv:2110.14168 (2021). https://arxiv.org/abs/2110.14168\\
{[}10{]} Huang et al.~\emph{RouterEval.} arXiv:2503.10657 (2025). https://arxiv.org/abs/2503.10657\\
{[}11{]} Qwen Team. \emph{Qwen3 Technical Report.} arXiv:2505.09388 (2025). https://arxiv.org/abs/2505.09388\\
{[}12{]} withmartian. \emph{RouterBench.} HuggingFace (2024). https://huggingface.co/datasets/withmartian/routerbench\\
{[}13{]} Salesforce AI Research. \emph{xRouter.} arXiv:2510.08439 (2025). https://arxiv.org/abs/2510.08439\\
{[}14{]} Chen et al.~\emph{RouterXBench.} arXiv:2602.11877 (2026). https://arxiv.org/abs/2602.11877\\
{[}15{]} Cassano et al.~\emph{MultiPL-E.} IEEE TSE (2022). arXiv:2208.08227. https://arxiv.org/abs/2208.08227\\
{[}16{]} Lihao Gu et al.~\emph{LLMRouterBench.} arXiv:2601.07206 (2025). https://arxiv.org/abs/2601.07206\\
{[}17{]} Liang et al.~\emph{HELM.} TMLR 2023. arXiv:2211.09110. https://arxiv.org/abs/2211.09110\\
{[}18{]} Zheng et al.~\emph{MT-Bench and Chatbot Arena.} NeurIPS 2023. arXiv:2306.05685. https://arxiv.org/abs/2306.05685\\
{[}19{]} Rai et al.~\emph{TaRo.} arXiv:2603.18411 (2025). https://arxiv.org/abs/2603.18411\\
{[}20{]} Tencent PlayCoder Team. \emph{PlayEval.} arXiv:2604.19742 (2026). https://arxiv.org/abs/2604.19742

\begin{center}\rule{0.5\linewidth}{0.5pt}\end{center}

\subsection{Limitations}\label{limitations}

\textbf{Single-user dataset.} All 15,363 queries come from one developer over 22 days. The DeBERTa classifier was trained and validated on the same distribution. Generalization to multi-user, domain-specific (legal, medical, scientific), or non-English deployments is not established. Classifier accuracy varies substantially by category: Reasoning 87.5\% (deterministic patterns dominate), Coding 46.3\% (harder to classify, more paraphrase variation). This imbalance will shift under different workloads. Single-developer training data risks encoding personal query patterns as universal signals --- the most significant limitation of this work.

\textbf{Benchmark sample sizes.} GPQA Diamond (20 questions) and GSM8K (20 problems) samples are too small for statistically robust conclusions; we report them as directional. Full MMLU (5,648 questions across 57 subjects) and HumanEval (164 problems, deterministic evaluation at temperature 0.0) have adequate statistical power. The GPQA 90\% result is interpretable only in light of the per-domain breakdown showing Tier 0 groundtruth resolution for definitional and computable questions.

\textbf{Hardware dependency.} Tier 2 (local inference) requires ≥24GB VRAM GPU. This is tested on an RTX 5090 32GB; Q4\_K\_M quantization narrows the requirement to \textasciitilde20GB. CPU-only deployments access Tiers 0+1 (38.9\% free), degrading to direct cloud routing for inference queries. 8GB VRAM deployments use a 9B-class model; the 100\% HumanEval result was achieved by a 9B model (OmniCoder-Qwen3.5-9B) and does not require the 30B tier.

\textbf{Privacy residual.} The 12.5\% of queries escalated to Tier 3 are transmitted to Anthropic, OpenAI, or Google servers under their data retention policies. No automatic PII detection or redaction is applied before escalation. Users with strict data-locality requirements should disable Tier 3 (\texttt{DISABLE\_CLOUD=1}).

\textbf{Evaluation methodology.} The 200-query human oracle evaluation was conducted by the same developer who generated the training data; it is not a double-blind randomized evaluation. RouterBench quality scores (75.0\%) use a threshold-based binary metric on a fixed 2024 dataset; this score is not directly comparable to HumanEval pass@1 or MMLU accuracy. The mµScore composite metric is author-defined and not yet externally validated.

\begin{center}\rule{0.5\linewidth}{0.5pt}\end{center}

\subsection{Ethical Considerations}\label{ethical-considerations}

\textbf{Data and privacy.} The 15,363-query evaluation dataset consists entirely of queries generated by the author during software development and research. No human subjects were recruited; no third-party user data was collected. Queries routed to cloud providers (12.5\%) were transmitted to Anthropic, OpenAI, and Google under their respective API terms; queries did not contain personally identifiable information beyond what is inherent in developer tool use.

\textbf{Environmental impact.} All local inference ran on a single workstation GPU (RTX 5090, \textasciitilde350W under load, Fedora Linux). Estimated local inference energy: \textasciitilde22 days × \textasciitilde8h/day active × 0.35kWh ≈ 62kWh. At 0.4 kg CO₂/kWh (US grid average), approximately 25 kg CO₂ equivalent. The 96.4\% cost reduction implies proportional reduction in cloud inference energy relative to an always-cloud baseline. A system that routes more queries to local inference reduces the aggregate energy demand on cloud data centers for equivalent query throughput.

\textbf{Dual use.} muLLM routes queries to capable language models; misuse of those models (harmful content generation, automated phishing) is equally possible through muLLM as through direct API access. muLLM does not add content moderation beyond provider-level policies. Operators should apply appropriate content policies at the application layer. The local-first architecture means content policy enforcement is the deploying organization's responsibility for Tier 2 queries.

\textbf{AI assistance.} Claude Code (Anthropic) assisted with software development and manuscript drafting throughout this project. Sessions used claude-sonnet-4-6 (primary), claude-opus-4-6/4-7 (review mode), and claude-haiku-4-5 (lightweight tasks). AI tools are not listed as authors per ACL policy. The routing system described here was itself used to route development queries --- a self-referential deployment that serves as an additional validity signal: a system too unreliable for sustained use could not have been built using itself.

\begin{center}\rule{0.5\linewidth}{0.5pt}\end{center}

\subsection{Appendix A: mμPETS --- Pareto Efficiency Timing Score}\label{appendix-a-mux3bcpets-pareto-efficiency-timing-score}

For routing systems, standard benchmarks (HumanEval, MMLU) measure the underlying model rather than the routing policy. mμPETS combines only the three dimensions where routing decisions have direct causal impact:

\begin{verbatim}
mμPETS = (routing_accuracy + cost_efficiency + ttft_score) / 3
\end{verbatim}

{\def\LTcaptype{none} % do not increment counter
\begin{longtable}[]{@{}llll@{}}
\toprule\noalign{}
Component & muLLM & Always-GPT-4 & Cloud-cheap only \\
\midrule\noalign{}
\endhead
\bottomrule\noalign{}
\endlastfoot
Routing accuracy & 0.824 & 1.00 & 0.644† \\
Cost efficiency & 0.987 & 0.00 & \textasciitilde0.30 \\
TTFT score & 0.85 & 0.30 & 0.35 \\
\textbf{mμPETS} & \textbf{0.887} & \textbf{0.43} & \textbf{\textasciitilde0.43} \\
\end{longtable}
}

\emph{†Cloud-cheap local win rate from MultiPL-E comparison (64.4\% vs 74.4\% local)}

A system routing everything to GPT-4 scores 0.0 on cost efficiency regardless of quality. mμPETS captures this directly. Both always-GPT-4 and always-cheap-cloud score \textasciitilde0.43; muLLM's 0.887 reflects that the routing policy itself adds substantial value beyond any single-model baseline.

\subsection{Appendix B: Controlled A/B Comparison}\label{appendix-b-controlled-ab-comparison}

100 held-out production queries routed through (A) full muLLM pipeline and (B) direct Sonnet API:

{\def\LTcaptype{none} % do not increment counter
\begin{longtable}[]{@{}lll@{}}
\toprule\noalign{}
Metric & muLLM & Sonnet direct \\
\midrule\noalign{}
\endhead
\bottomrule\noalign{}
\endlastfoot
Total cost & \textbf{\$0.0025} & \$0.5627 \\
Mean latency & \textbf{5,027ms} & 6,481ms \\
Cache hit rate & \textbf{41\%} & 0\% \\
Local tier rate & \textbf{57\%} & 0\% \\
\end{longtable}
}

99.6\% cost reduction, 22\% lower latency, on held-out production queries with already-warmed cache.
